\section{Results}
\label{sec:results}

\subsection{Prompt Format Performance}
\label{sec:format_performance}

\tabref{tab:formats} shows the accuracy of each prompt format on \ssttwo.
The best format (\texttt{caps\_instruction}) achieves 71.0\% accuracy, while the worst formats (\texttt{minimal}, \texttt{emoji}, \texttt{bare\_newline}) score at chance (50.0\%), yielding a 21 percentage point spread.
This confirms that prompt sensitivity exists even on a simple binary classification task with a small model, consistent with prior findings \citep{sclar2023quantifying}.

\begin{table}[t]
    \centering
    \caption{Accuracy of 20 prompt formats on \ssttwo with \gptsmall (200 examples). Formats above the median (0.603) are labeled ``Stable''; those below are labeled ``Failure.'' Best result in \textbf{bold}.}
    \label{tab:formats}
    \begin{tabular}{@{}llc@{}}
        \toprule
        \textbf{Format} & \textbf{Accuracy (\%)} & \textbf{Category} \\
        \midrule
        \texttt{caps\_instruction} & \textbf{71.0} & Stable \\
        \texttt{xml\_tags}         & \textbf{71.0} & Stable \\
        \texttt{academic}          & 68.5 & Stable \\
        \texttt{chat}              & 68.5 & Stable \\
        \texttt{classify}          & 68.0 & Stable \\
        \texttt{reverse\_order}    & 68.0 & Stable \\
        \texttt{formal}            & 65.0 & Stable \\
        \texttt{parenthetical}     & 64.5 & Stable \\
        \texttt{simple\_direct}    & 61.5 & Stable \\
        \texttt{pipe\_sep}         & 60.5 & Stable \\
        \cmidrule(lr){1-3}
        \texttt{dash\_sep}         & 60.0 & Failure \\
        \texttt{colon\_nospace}    & 60.0 & Failure \\
        \texttt{question}          & 59.0 & Failure \\
        \texttt{star\_rating}      & 55.0 & Failure \\
        \texttt{numbered}          & 52.5 & Failure \\
        \texttt{json\_like}        & 52.5 & Failure \\
        \texttt{tab\_sep}          & 52.0 & Failure \\
        \texttt{minimal}           & 50.0 & Failure \\
        \texttt{emoji}             & 50.0 & Failure \\
        \texttt{bare\_newline}     & 50.0 & Failure \\
        \bottomrule
    \end{tabular}
\end{table}

The stable formats share common properties: they provide clear task instructions, use standard label words (positive/negative), and structurally delineate the input from the instruction (XML tags, all-caps headers, explicit classification commands).
The failure formats either provide minimal context (\texttt{bare\_newline}, \texttt{minimal}), use non-standard label tokens (\texttt{emoji}, \texttt{star\_rating}), or lack clear task boundaries (\texttt{tab\_sep}, \texttt{json\_like}).

\subsection{Geometric Feature Correlations}
\label{sec:correlations}

\tabref{tab:correlations} presents the Spearman rank correlations between the top geometric features and prompt accuracy.
Six of 23 features achieve statistical significance at $p < 0.05$.

\begin{table}[t]
    \centering
    \caption{Top geometric features ranked by absolute Spearman correlation with prompt accuracy. Asterisks (*) denote $p < 0.05$. Bootstrap 95\% CIs computed with 10{,}000 resamples. The permutation test for the top feature yields $p = 0.036$ (5{,}000 permutations).}
    \label{tab:correlations}
    \resizebox{\textwidth}{!}{%
    \begin{tabular}{@{}lccc@{}}
        \toprule
        \textbf{Feature} & \textbf{Spearman} $\spearman$ & \textbf{$p$-value} & \textbf{95\% Bootstrap CI} \\
        \midrule
        Mean Frobenius norm       & $+0.474$ & $0.035^*$ & $[-0.004, +0.808]$ \\
        Variance layer std        & $-0.472$ & $0.036^*$ & $[-0.815, -0.014]$ \\
        Mean attention variance   & $-0.460$ & $0.041^*$ & $[-0.793, +0.004]$ \\
        Std effective rank        & $+0.454$ & $0.044^*$ & $[-0.024, +0.791]$ \\
        Late-layer entropy        & $+0.448$ & $0.048^*$ & $[-0.008, +0.786]$ \\
        Mean SV entropy           & $+0.446$ & $0.049^*$ & --- \\
        \midrule
        Mean entropy              & $+0.412$ & $0.071$ & --- \\
        Mid-layer variance        & $-0.389$ & $0.090$ & --- \\
        Std Frobenius norm        & $+0.345$ & $0.137$ & --- \\
        \bottomrule
    \end{tabular}
    }
\end{table}

The directions of the significant correlations are interpretable:
\begin{itemize}[leftmargin=*,itemsep=0pt,topsep=0pt]
    \item \textbf{Mean Frobenius norm} ($\spearman = +0.474$): Higher overall attention magnitude correlates with better performance. Stronger attention signals may indicate that the model can reliably extract task-relevant information from the prompt.
    \item \textbf{Variance layer std} ($\spearman = -0.472$): Greater inconsistency of attention variance across layers correlates with worse performance. This aligns with the hypothesis that geometric instability---measured as cross-layer inconsistency---predicts failure.
    \item \textbf{Late-layer entropy} ($\spearman = +0.448$): Higher entropy in the final layers (9--12) correlates with better performance. More distributed attention in late layers may reflect richer integration of contextual information.
\end{itemize}

\para{Permutation test.}
To control for multiple comparisons across 23 features, we run a permutation test for the top-ranked feature (mean Frobenius norm).
We permute the accuracy labels 5{,}000 times and compute the maximum absolute Spearman correlation across all features in each permutation.
The observed correlation of $|\spearman| = 0.474$ exceeds 96.4\% of permuted maxima, yielding a permutation $p = 0.036$.
This confirms that the top correlation is unlikely to be spurious.

\subsection{Failure Classification}
\label{sec:classification}

\tabref{tab:classifiers} reports leave-one-out classification results.

\begin{table}[t]
    \centering
    \caption{Binary failure classification results (LOO cross-validation, $n = 20$). The logit-confidence baseline always predicts one class, scoring 0\% on the minority class. Best geometry-based result in \textbf{bold}.}
    \label{tab:classifiers}
    \begin{tabular}{@{}lcccc@{}}
        \toprule
        \textbf{Method} & \textbf{Accuracy (\%)} & \textbf{AUC} & \textbf{Precision} & \textbf{F1} \\
        \midrule
        Random baseline       & 50.0 & 0.500 & 0.500 & 0.500 \\
        Logit confidence      & 0.0  & 0.000 & 0.000 & 0.000 \\
        \midrule
        Top-2 features (LR)   & 55.0 & 0.470 & 0.545 & 0.571 \\
        Top-3 features (LR)   & \textbf{60.0} & \textbf{0.520} & \textbf{0.600} & \textbf{0.600} \\
        Top-5 features (LR)   & \textbf{60.0} & 0.480 & \textbf{0.600} & \textbf{0.600} \\
        \bottomrule
    \end{tabular}
\end{table}

The logit-confidence baseline completely fails, always predicting a single class.
This indicates that the maximum softmax probability does not distinguish good from bad prompt formats---a surprising result that strongly motivates looking beyond output-based heuristics.

The geometry-based classifiers achieve 60\% LOO accuracy with the top-3 and top-5 features, a 10 percentage point improvement over random guessing.
While this improvement is modest, the small sample size ($n = 20$) severely limits classifier generalization.
Adding more features beyond the top 3 does not improve performance, suggesting that the marginal features add noise rather than signal at this sample size.
